\chapter{Calibration Controller}
\label{CalContr}

This chapter provides guidelines about the usage of the APIs configured in the Calibration Controller.

\section{Initialization}
Please remember to connect the client-server interface SrvcMgr\_CalibrationController to properly allow the calibration of the services and service groups.

\section{Runnable\_SrvcMgr\_DisableService}
SUM\_SrvcMgr implements the Runnable\_SrvcMgr\_DisableService that will be called through the C/S interface SrvcMgr\_ServiceCalibrationOff (operation SrvcMgr\_DisableService) 
by the SWCs to change the state (disable) of services through SD APIs.

This Runnable take as input parameter the ID value of the Service that the user want to disable. 

Please, take into account that ID refer to the SdClientServiceId or SdServerServiceId specified in the container Sd/SdConfig/SdInstance/SdClientService for a Client Service or 
Sd/SdConfig/SdInstance/SdServerService for a Service Service. The Runnable Runnable\_SrvcMgr\_DisableService understand by itself if the ID received correspond to a Client or Server Service.

This Runnable will change the status of a Service through SD APIs:

\begin{itemize}
\item Call to Sd\_ServerServiceSetState()/Sd\_ClientServiceSetState();
\item Call to Sd\_CalibrateServicesOff();
\end{itemize}

The Runnable\_SrvcMgr\_DisableService return the following value: 
\begin{itemize}
\item E\_OK: if the Service is properly disabled.
\item E\_NOT\_OK: if the Service is NOT properly disabled for the following reason:
\begin{itemize}
\item The Service ID is not configured in the SD container; 
\item One of the SD APIs called return E\_NOT\_OK;
\end{itemize}
\end{itemize}

Please, refer to the SD module specification in order to understand the meaning of possible E\_NOT\_OK received by the SD APIs. 
The user SHALL connect the port PP\_SrvcMgr\_ServiceCalibrationOff to their SWC and configure the call in synchronous way in order to trigger this Runnable properly.

\section{Runnable\_SrvcMgr\_DisableServiceGroup}
SUM\_SrvcMgr implements the Runnable\_SrvcMgr\_DisableServiceGroup that will be called through the C/S interface SrvcMgr\_ServiceCalibrationOff (operation SrvcMgr\_DisableServiceGroup) 
by the SWCs to change the state (disable) of services group through SD APIs.

This Runnable takes, as an input parameter, the ID value of the Group Service that the user wants to disable. 

Please, take into account that ID refer to the SdServiceGroupHandleId specified in the container Sd/SdConfig/SdServiceGroup.

This Runnable will change the status of a Group Service through SD APIs:

\begin{itemize}
\item Call to Sd\_ServiceGroupStop();
\item Call to Sd\_CalibrateServicesOff();
\end{itemize}

The Runnable\_SrvcMgr\_DisableServiceGroup return the following value: 
\begin{itemize}
\item E\_OK: The Service Group is properly disabled.
\item E\_NOT\_OK: if the Service Group is NOT properly disabled for the following reason:
\begin{itemize}
\item There is no Service Group configured in the SD container; 
\item The Service Group ID is not configured in the SD container; 
\item The Service Group does not have at least one service associated; 
\item The API Sd\_CalibrateServicesOff() returns E\_NOT\_OK;
\end{itemize}
\end{itemize}

Please, refer to the SD module specification in order to understand the meaning of possible E\_NOT\_OK received by the SD API.\\
The user SHALL connect the port PP\_SrvcMgr\_ServiceCalibrationOff to their SWC and configure the call in synchronous way in order to trigger this Runnable properly.\\
Moreover, the integrator SHOULD follow the normal workflow about the generation: firstly generate the BSW and then the SrvcMgr. This is necessary since the Sd/SdConfig/SdServiceGroup/SdServiceGroupHandleId parameter
is a forwarded value: the integrator SHALL not put a value there because it is already done during the BSW generation by itself. 